\documentclass[a4paper, 12pt]{article}
% math symbols
\usepackage{amssymb}
\usepackage{amsmath}
\usepackage{mathrsfs}
\usepackage{mathseries}


\usepackage[margin = 2cm]{geometry}

\tolerance = 1000
\emergencystretch = 0.74cm



\pagestyle{empty}
\parindent = 0mm

\renewcommand{\coursetitle}{DM/ML}
\setcounter{curtask}{1}

\setmathstyle{\includegraphics[scale = 0.05]{../pics/utia-rest.png}}{Мат. логика}{3 курс}


\begin{document}

\begin{definition*}
    Пусть $A$~--- $\sigma$-структура, $X \subseteq A$, и $\sigma_x \coloneqq (\sigma, \{b^x \mid x \in
    X\})$~--- расширение сигнатуры $\sigma$ при помощи новых констант для каждого элемента множества
    $X$. $A_X$, соответственно, расширение $A$, где элементы $b^x$ интерпретируются естественным образом
    (т.е для любого $a \in X$ константа $b^a$ интерпретируется, как $a$).

    \deftext{Частичным $n$-типом} (или просто $n$-типом) $\sigma$-структуры $A$ над множеством $X$
    называть типы теории $\Theory(A_x)$, соответствющим образом перенесём понятия \deftext{полного типа}
    и \deftext{реализации типа}. По умолчанию будем считать множество $X$ пустым.
\end{definition*}

\task{
    Предъявите пример сигнатуры $\sigma$ и такой её интерпретации $M$, что найдется полный $1$-тип,
    который не реализуется в $M$.
}


\libproblem{math-logic}{countable-model-uncount-types}

\task{
    Пусть $M$~--- $\sigma$-структура и $p$ некоторый её $n$-тип. Покажите, что существует такое
    элементарное расширение структуры $M$, в котором $p$ реализуется.
}

\task{
    Пусть $\sigma$~--- счетная сигнатура и $T$~--- полная $\sigma$-теория, у которой существует
    бесконечная модель. Покажите, что если $T$~--- $\aleph_0$-категорична, то все $n$-типы реализуются в
    счетной модели данной теории.
}


\task{
    Рассмотрим теорию $\Theory(\mathbb{Z}, s)$, $s(x) \coloneqq x + 1$. Опишите полные $n$-типы.
}

\task{
    Рассмотрим теорию $\Theory(\mathbb{Z}, s, <)$, $s(x) \coloneqq x + 1$. Опишите полные $n$-типы.
}




\end{document}



%%% Local Variables:
%%% mode: latex
%%% TeX-master: t
%%% End: